\subsection{Cohomological Description}

Let G be a group, let F be an abelian group, and let \(\tau: G \to \text{Aut}(F)\) be a group homomorphism. For any \(g \in G\) and \(f \in F\) , we also write \(^{g}f\) for \(\tau(g) \cdot f\) . A 2-cocycle of G with values in F is a mapping \(c: G \times G \to F\) such that for every \((g_{1}, g_{2}, g_{3}) \in G \times G \times G\) , we have

\[
{ } ^ { g _ { 1 } } c ( g _ { 2 } , g _ { 3 } ) c ( g _ { 1 } , g _ { 2 } g _ { 3 } ) = c ( g _ { 1 } , g _ { 2 } ) c ( g _ { 1 } g _ { 2 } , g _ { 3 } ) .\tag{7}
\]

Since F is commutative, the set of 2-cocycles is a subgroup of the group of mappings from \(G \times G\) to F; we denote it by \(Z^{2}(G,F)\) . We denote the group of mappings from G to F by \(C^{1}(G,F)\) . For every \(h \in C^{1}(G,F)\) and \((g_{1}, g_{2}) \in G \times G\) , set

\[
(\partial h) (g _ {1}, g _ {2}) = ^ {g _ {1}} h (g _ {2}) h (g _ {1} g _ {2}) ^ {- 1} h (g _ {1}).\tag{8}
\]

A simple calculation shows that the mapping \(\partial h: G \times G \to F\) is a 2-cocycle. The mapping \(\partial: C^{1}(G, F) \to Z^{2}(G, F)\) is a group homomorphism. For any \(h \in C^{1}(G, F)\) , the 2-cocycle \(\partial h\) is called a 2-coboundary. The group

\(Z^{2}(G,F)/\partial(C^{1}(G,F))\) is denoted by \(H^{2}(G,F)\) and is called the second cohomology group of G with coefficients in F.

*The notation above agrees with that of X, §6, n° 8, p.~112 concerning group cohomology.*

Let \(\mathcal{E} = (\Gamma, \iota, \pi)\) be a \(\tau\) -extension. Let \(\sigma\) be a section of the surjective mapping \(\pi\) (Set Theory, II, §3, No.~8, p.~86), that is, a mapping from G to \(\Gamma\) such that \(\pi(\sigma(g)) = g\) for every \(g\) in G. For every \(g_1, g_2 \in G\) , \(\sigma(g_1)\sigma(g_2)\sigma(g_1g_2)^{-1}\) belongs to \(\mathrm{Ker}(\pi)\) , so that there exists a unique mapping \(c_\sigma: G \times G \to F\) such that

\[
\iota (c _ {\sigma} (g _ {1}, g _ {2})) = \sigma (g _ {1}) \sigma (g _ {2}) \sigma (g _ {1} g _ {2}) ^ {- 1}\tag{9}
\]

for all \(g_{1}, g_{2} \in G\) . The mapping \(c_{\sigma}\) is constant with value 1 if and only if \(\sigma\) is a group homomorphism.

PROPOSITION 6. --- The mapping \(c_{\sigma}\) is an element of the group \(Z^{2}(G,F)\) , and its class in the cohomology group \(H^{2}(G,F)\) depends only on the class of the \(\tau\) -extension E.

We say that the mapping \(c_{\sigma}\) is the 2-cocycle associated with \(\sigma\) and that its cohomology class in \(\mathrm{H}^2 (\mathrm{G},\mathrm{F})\) is the cohomology class of the \(\tau\) -extension \(\mathcal{E}\) .

Let us first verify that \(c_{\sigma}\) satisfies the cocycle condition. Let \(g_{1}, g_{2}\) , and \(g_{3}\) be elements of G. Using formula (1) of VIII, p.~285 and (9), we obtain the relations

\[
\begin{array}{r l} & {\iota (g _ {1} c _ {\sigma} (g _ {2}, g _ {3}) c _ {\sigma} (g _ {1}, g _ {2} g _ {3}))} \\ & {\quad = \sigma (g _ {1}) \sigma (g _ {2}) \sigma (g _ {3}) \sigma (g _ {2} g _ {3}) ^ {- 1} \sigma (g _ {1}) ^ {- 1} \sigma (g _ {1}) \sigma (g _ {2} g _ {3}) \sigma (g _ {1} g _ {2} g _ {3}) ^ {- 1}} \\ & {\quad = \sigma (g _ {1}) \sigma (g _ {2}) \sigma (g _ {3}) \sigma (g _ {1} g _ {2} g _ {3}) ^ {- 1}} \end{array}
\]

and

\[
\begin{array}{r l} & {\iota (c _ {\sigma} (g _ {1}, g _ {2}) c _ {\sigma} (g _ {1} g _ {2}, g _ {3}))} \\ & {\qquad = \sigma (g _ {1}) \sigma (g _ {2}) \sigma (g _ {1} g _ {2}) ^ {- 1} \sigma (g _ {1} g _ {2}) \sigma (g _ {3}) \sigma (g _ {1} g _ {2} g _ {3}) ^ {- 1}} \\ & {\qquad = \sigma (g _ {1}) \sigma (g _ {2}) \sigma (g _ {3}) \sigma (g _ {1} g _ {2} g _ {3}) ^ {- 1}.} \end{array}
\]

The mapping \(c_{\sigma}\) is therefore an element of \(\mathbf{Z}^2 (\mathrm{G},\mathrm{F})\) .

Let us now prove that the image of \(c_{\sigma}\) in \(\mathrm{H}^{2}(\mathrm{G},\mathrm{F})\) is independent of the chosen section \(\sigma\) . Let \(\sigma\) and \(\sigma'\) be such sections. For every element g of G, there exists a unique element \(a_{g}\) of F such that \(\sigma'(g) = \iota(a_{g})\sigma(g)\) . Let \(g_{1}\) and \(g_{2}\) be elements of G. Using the definition (9), we obtain the equalities

\[
\begin{array}{r l} & {\iota (c _ {\sigma^ {\prime}} (g _ {1}, g _ {2})) = \sigma^ {\prime} (g _ {1}) \sigma^ {\prime} (g _ {2}) \sigma^ {\prime} (g _ {1} g _ {2}) ^ {- 1}} \\ & {\qquad = \iota (a _ {g _ {1}}) \sigma (g _ {1}) \iota (a _ {g _ {2}}) \sigma (g _ {2}) \sigma (g _ {1} g _ {2}) ^ {- 1} \iota (a _ {g _ {1} g _ {2}}) ^ {- 1}} \\ & {\qquad = \iota (a _ {g _ {1}}) \iota (^ {g _ {1}} a _ {g _ {2}}) \iota (c _ {\sigma} (g _ {1}, g _ {2})) \iota (a _ {g _ {1} g _ {2}}) ^ {- 1}.} \end{array}\tag{10}
\]

Since the group F is commutative, we have the relations

\[
c _ {\sigma^ {\prime}} (g _ {1}, g _ {2}) c _ {\sigma} (g _ {1}, g _ {2}) ^ {- 1} = (^ {g _ {1}} a _ {g _ {2}}) a _ {g _ {1} g _ {2}} ^ {- 1} a _ {g _ {1}} = (\partial a) (g _ {1}, g _ {2}),\tag{11}
\]

as follows from (8). This proves that the classes of \(c_{\sigma'}\) and \(c_{\sigma}\) in \(\mathrm{H}^2(\mathrm{G},\mathrm{F})\) coincide.

Finally, let \(\mathcal{E} = (\Gamma, \iota, \pi)\) and \(\mathcal{E}' = (\Gamma', \iota', \pi')\) be isomorphic \(\tau\) -extensions, let \(u: \mathcal{E} \to \mathcal{E}'\) be a morphism of \(\tau\) -extensions, and let \(\sigma\) be a section of the mapping \(\pi\) . The mapping \(u \circ \sigma\) is a section of the mapping \(\pi'\) , and we have \(c_{\sigma} = c_{u \circ \sigma}\) . The class of \(c_{\sigma}\) in \(\mathrm{H}^2(\mathrm{G}, \mathrm{F})\) therefore depends only on the class of \(\mathcal{E}\) in \(\mathrm{Ex}_{\tau}(\mathrm{G}, \mathrm{F})\) .

We denote by \(\Theta_{\tau}: \operatorname{Ex}_{\tau}(\mathrm{G},\mathrm{F}) \to \mathrm{H}^{2}(\mathrm{G},\mathrm{F})\) , or more simply \(\Theta\) , the mapping that sends the class of a \(\tau\) -extension \((\Gamma,\iota,\pi)\) to the class of the 2-cocycle \(c_{\sigma}\) , where \(\sigma\) is a section of the surjection \(\pi\) .

THEOREM 1. --- The mapping \(\Theta\) is a group isomorphism from \(\operatorname{Ex}_{\tau}(\mathrm{G},\mathrm{F})\) to \(\mathrm{H}^2 (\mathrm{G},\mathrm{F})\) .

Let us first prove that \(\Theta\) is a group homomorphism. Let \(\mathcal{E} = (\Gamma, \iota, \pi)\) and \(\mathcal{E}' = (\Gamma', \iota', \pi')\) be \(\tau\) -extensions, and let \(\sigma\) (resp. \(\sigma'\) ) be a section of the mapping \(\pi\) (resp. \(\pi'\) ). We denote by \(\mathcal{E}\mathcal{E}' = (\Gamma'', \iota'', \pi'')\) the product of the \(\tau\) -extensions E and \(E'\) . We denote by \([\gamma, \gamma']\) the image in \(\Gamma''\) of an element \((\gamma, \gamma')\) of \(\Gamma \times_{G} \Gamma'\) by the surjective homomorphism from the remark of VIII, p.~293. The mapping from G to \(\Gamma''\) that sends an element g to \([\sigma(g), \sigma'(g)]\) is a section \(\sigma''\) of the mapping \(\pi''\) . Let \(g_1\) and \(g_2\) be elements of G. We have the relations

\[
\begin{array}{r l} & {\iota^ {\prime \prime} (c _ {\sigma^ {\prime \prime}} (g _ {1}, g _ {2})) = \sigma^ {\prime \prime} (g _ {1}) \sigma^ {\prime \prime} (g _ {2}) \sigma^ {\prime \prime} (g _ {1} g _ {2}) ^ {- 1}} \\ & {\qquad = [ \iota (c _ {\sigma} (g _ {1}, g _ {2})), \iota^ {\prime} (c _ {\sigma^ {\prime}} (g _ {1}, g _ {2}) ]} \\ & {\qquad = \iota^ {\prime \prime} (c _ {\sigma} (g _ {1}, g _ {2}) c _ {\sigma^ {\prime}} (g _ {1}, g _ {2})).} \end{array}
\]

We therefore have \(c_{\sigma''}(g_1, g_2) = c_\sigma(g_1, g_2)c_{\sigma'}(g_1, g_2)\) .

Let us prove that the mapping \(\Theta\) is injective. Let \(\mathcal{E} = (\Gamma, \iota, \pi)\) be a \(\tau\) -extension, and let \(\sigma\) be a section of the mapping \(\pi\) such that the image of \(c_{\sigma}\) in \(\mathrm{H}^{2}(\mathrm{G}, \mathrm{F})\) is trivial. There exists a mapping \(a : G \to F\) such that

\[
c _ {\sigma} (g _ {1}, g _ {2}) = (^ {g _ {1}} a (g _ {2})) a (g _ {1} g _ {2}) ^ {- 1} a (g _ {1})
\]

for all \(g_{1}, g_{2} \in G\) . We then define \(\sigma' : G \to \Gamma\) by

\[
\sigma^ {\prime} (g) = \iota (a (g) ^ {- 1}) \sigma (g)
\]

for \(g \in G\) . Using (11), we obtain that \(c_{\sigma'}\) is constant with value 1; consequently, \(\sigma'\) is a group homomorphism, which proves that the \(\tau\) -extension E is semitrivial (I, §6, No.~1, p.~67, Proposition 4).

Let us prove that the mapping \(\Theta\) is surjective. Let c be an element of \(Z^{2}(G,F)\) . We endow the set \(\Gamma = F \times G\) with the following law of composition:

\[
(f _ {1}, g _ {1}) (f _ {2}, g _ {2}) = \left(f _ {1} (^ {g _ {1}} f _ {2}) c (g _ {1}, g _ {2}), g _ {1} g _ {2}\right)\tag{12}
\]

for all \(f_{1}, f_{2} \in F\) and \(g_{1}, g_{2} \in G\) . We have the relations

\[
\begin{array}{r l} (f _ {1}, g _ {1}) ((f _ {2}, g _ {2}) (f _ {3}, g _ {3})) & = (f _ {1}, g _ {1}) \big (f _ {2} (^ {g _ {2}} f _ {3}) c (g _ {2}, g _ {3}), g _ {2} g _ {3} \big) \\ & = \big (f _ {1} (^ {g _ {1}} f _ {2}) (^ {g _ {1} g _ {2}} f _ {3}) (^ {g _ {1}} c (g _ {2}, g _ {3})) c (g _ {1}, g _ {2} g _ {3}), g _ {1} g _ {2} g _ {3} \big) \end{array}
\]

and

\[
\begin{array}{c} ((f _ {1}, g _ {1}) (f _ {2}, g _ {2}))   (f _ {3}, g _ {3}) = \big (f _ {1} (^ {g _ {1}} f _ {2}) c (g _ {1}, g _ {2}), g _ {1} g _ {2} \big) (f _ {3}, g _ {3}) \\ = \big (f _ {1} (^ {g _ {1}} f _ {2}) (^ {g _ {1} g _ {2}} f _ {3}) c (g _ {1}, g _ {2}) c (g _ {1} g _ {2}, g _ {3}), g _ {1} g _ {2} g _ {3} \big) \end{array}
\]

for all \(f_{1}, f_{2}, f_{3} \in F\) and \(g_{1}, g_{2}, g_{3} \in G\) . It follows that this law is associative because c is a 2-cocycle. For every \((f, g) \in \Gamma\) , we have

\[
(f, g) (c (e, e) ^ {- 1}, e) = (f (^ {g} c (e, e) ^ {- 1}) c (g, e), g).
\]

Now, it follows from the definition of a 2-cocycle that \(c(g, e) = {}^g c(e, e)\) , from which we deduce that \((f, g)(c(e, e)^{-1}, e) = (f, g)\) . We prove likewise that

\[
(c (e, e) ^ {- 1}, e) (f, g) = (f c (e, e) ^ {- 1} c (e, g), g) = (f, g).
\]

The law of composition defined by (12) therefore has identity element \((c(e,e)^{-1},e)\) , and every element \((f,g)\) of \(\Gamma\) is invertible, with inverse

\[
\left(^ {g ^ {- 1}} (f ^ {- 1} c (e, e) ^ {- 1} c (g, g ^ {- 1}) ^ {- 1}), g ^ {- 1}\right).
\]

So we have endowed \(\Gamma\) with a group structure. Let us denote by \(\iota: F \to G\) the mapping that sends f to \((c(e, e)^{-1}f, e)\) , by \(\pi: \Gamma \to G\) the second projection, and by \(\sigma\) the mapping \(g \mapsto (1, g)\) . The mappings \(\iota\) and \(\pi\) are group homomorphisms, so the triple \(\mathcal{E} = (\Gamma, \iota, \pi)\) is a \(\tau\) -extension, \(\sigma\) is a section of the mapping \(\pi\) , and the associated cocycle \(c_{\sigma}\) is equal to c because

\[
\begin{array}{c} \sigma (g _ {1}) \sigma (g _ {2}) \sigma (g _ {1} g _ {2}) ^ {- 1} = (1, g _ {1}) (1, g _ {2}) (1, g _ {1} g _ {2}) ^ {- 1} = (c (g _ {1}, g _ {2}), g _ {1} g _ {2}) (1, g _ {1} g _ {2}) ^ {- 1} \\ = (c (g _ {1}, g _ {2}) c (e, g _ {1} g _ {2}) ^ {- 1}, e) = (c (e, e) ^ {- 1} c (g _ {1}, g _ {2}), e) \end{array}
\]

for \(g_{1}, g_{2} \in G\) .

Remark. --- Let G be a group, let F and \(F'\) be abelian groups, let \(\tau\) (resp. \(\tau'\) ) be a group homomorphism from G to the automorphism group of F (resp. \(F'\) ), and let \(v : F \to F'\) be a group morphism such that

\[
v (\tau (g) \cdot f) = \tau^ {\prime} (g) \cdot v (f)\tag{13}
\]

for \(g \in G\) and \(f \in F\) . Let \(\alpha\) be an element of \(\operatorname{Ex}_{\tau}(\mathrm{G}, \mathrm{F})\) . If the cocycle c represents \(\Theta(\alpha)\) , then \(v \circ c\) represents \(\Theta(v_{*}(\alpha)) \in \mathrm{H}^{2}(\mathrm{G}, \mathrm{F}')\) .

